\section{Experiments II: Logical Propositions}
\label{sec:experiments-logic}

The experiments in this section instantiate the MDL objective function
in the domain of logical proposition datasets, using the QBBN coding
scheme developed in Section~\ref{sec:coding}. All results are
reproducible from the \texttt{mdl-science} repository.

A dataset is a list of boolean proposition observations. A theory is a
set of \emph{implication links} (premise $\to$ conclusion, firing when
the premise is true) and \emph{conjunctive links}
($[\text{premise}_1, \ldots, \text{premise}_n] \to$ conclusion, firing
only when all premises are true). These correspond directly to the QBBN:
conjunctive links are AND gates, implication links are OR-gate inputs.

The two-part code is concrete. $L(H)$ counts the number of links times
the bits per weight---a fixed cost per link, independent of dataset
size. $L(D \mid H)$ is the total negative log-likelihood of the observed
propositions under the theory's conditional probabilities, with link
weights calibrated empirically from the data. Propositions not covered
by any link fall back to their empirical base rate, ensuring the theory
can only help---never hurt---on uncovered propositions.

\subsection{Conjunctive Structure Wins (\texttt{mdl compress})}

Data where $r = p \wedge q$ exactly ($n = 50$ triples). Six theories
compete. Results are reported as MDL score and bits saved versus null
(positive = beats no theory; negative = worse than no theory):

\begin{center}
\begin{tabular}{lrrrr}
\toprule
Theory & $L(H)$ & $L(D \mid H)$ & MDL & Bits saved \\
\midrule
null (no links)                 &  0.0 & 143.5 & 143.5 &      $0.0$ \\
one link: $p \to r$             &  8.0 & 134.2 & 142.2 &    $+1.3$ \\
one link: $q \to r$             &  8.0 & 140.0 & 148.0 &    $-4.5$ \\
two links: $p \to r$, $q \to r$ & 16.0 & 141.8 & 157.8 &   $-14.3$ \\
conjunctive: $p \wedge q \to r$ &  8.0 & 117.1 & 125.1 & $+\mathbf{18.4}$ \\
wrong link: $p \to q$           &  8.0 & 143.2 & 151.2 &    $-7.7$ \\
\bottomrule
\end{tabular}
\end{center}

\textit{Bits saved = MDL(null) $-$ MDL(theory).}

The conjunctive theory wins by 18.4 bits using only 8 bits to describe
one link. The disjunctive two-link theory loses by 14.3 bits despite
having both premises available: it misfires when only one of $p$ or $q$
is true, predicting $r = 1$ when $r = 0$. The wrong link loses by 7.7
bits. Even the partial theory $p \to r$ barely wins at $+1.3$ bits.

The key result: having more links is not always better. The two-link
disjunctive theory has strictly more information about $r$ than the
one-link conjunctive theory, yet it performs worse. What matters is
whether the links capture the true generative structure. The disjunctive
theory has the right variables but the wrong logical connective.

This result is a direct consequence of the QBBN coding scheme. The AND
gate (conjunctive link) is a zero-parameter, deterministic structure
that adds only its own description to $L(H)$. The OR-gate inputs
(implication links) each add weight parameters. When the true structure
is conjunctive, the conjunctive coding is maximally efficient: it pays
zero for the AND computation itself, only for the single weight that
says how strongly the conjunction predicts the conclusion.

\subsection{Theories Are Amortized Over Data (\texttt{mdl scale})}

One conjunctive link, $L(H) = 8$ bits, regularity $= 1.0$, $n$ from
5 to 1000:

\begin{center}
\begin{tabular}{rrrrr}
\toprule
$n$ & $L(D)$ & MDL & Bits saved & Comp.\% \\
\midrule
5    &   13.3 &   18.7 &   $-5.4$ & $-40.2$\% \\
10   &   28.8 &   31.6 &   $-2.8$ &  $-9.7$\% \\
25   &   70.9 &   65.8 &   $+5.0$ &  $+7.1$\% \\
50   &  136.7 &  118.1 &  $+18.5$ & $+13.6$\% \\
100  &  276.2 &  231.2 &  $+44.9$ & $+16.3$\% \\
200  &  535.7 &  440.5 &  $+95.2$ & $+17.8$\% \\
500  & 1326.8 & 1079.0 & $+247.9$ & $+18.7$\% \\
1000 & 2680.7 & 2170.1 & $+510.6$ & $+19.0$\% \\
\bottomrule
\end{tabular}
\end{center}

\textit{Bits saved = $L(D) -$ MDL. Positive means the theory beats the
no-theory baseline.}

$L(H) = 8$ bits at every scale. The theory loses at $n = 5$ and
$n = 10$, crosses into positive territory between $n = 10$ and $n = 25$,
and saves 510 bits at $n = 1000$---64 times its own description length.

This is the formal version of the intuition behind scientific laws.
Newton's laws cost a fixed number of bits to state; every new
observation is compressed for free once the theory exists. The
amortization is not a metaphor. It is a computable quantity,
demonstrated here at the level of individual bits, with the crossover
point (theory overhead recovered at $\approx n = 15$ pairs) precisely
measurable.

The QBBN coding scheme makes this particularly clean. Because $L(H)$
counts only the structural description of the links (which propositions,
which logical connective, what weight precision), it is genuinely
independent of dataset size. The same 8 bits describes the theory
whether the dataset has 5 observations or 5 million. This is the coding
scheme doing its job: separating the cost of the theory from the cost
of the data.

\subsection{MDL Selects Correct Theory Complexity (\texttt{mdl compete})}

Four theories at three regularity levels. T0 = null, T1 = one
implication link, T2 = correct conjunctive link, T3 = T2 plus one
spurious implication link:

\begin{center}
\begin{tabular}{lrrrr}
\toprule
Theory & $L(H)$ & $L(D \mid H)$ & MDL & Bits saved \\
\midrule
\multicolumn{5}{l}{\textit{regularity = 1.0}} \\
T0: null                &  0.0 & 276.2 & 276.2 &       $0.0$ \\
T1: one link            &  8.0 & 272.6 & 280.6 &      $-4.5$ \\
T2: correct conjunctive &  8.0 & 223.2 & 231.2 & $+\mathbf{44.9}$ \\
T3: T2 + spurious link  & 16.0 & 223.2 & 239.2 &     $+36.9$ \\
\midrule
\multicolumn{5}{l}{\textit{regularity = 0.5}} \\
T0: null                &  0.0 & 272.5 & 272.5 & $\mathbf{0.0}$ \\
T1: one link            &  8.0 & 272.0 & 280.0 &      $-7.6$ \\
T2: correct conjunctive &  8.0 & 262.6 & 270.6 &      $+1.9$ \\
T3: T2 + spurious link  & 16.0 & 262.6 & 278.6 &      $-6.1$ \\
\midrule
\multicolumn{5}{l}{\textit{regularity = 0.0}} \\
T0: null                &  0.0 & 273.5 & 273.5 & $\mathbf{0.0}$ \\
T1: one link            &  8.0 & 273.5 & 281.5 &      $-8.0$ \\
T2: correct conjunctive &  8.0 & 273.5 & 281.5 &      $-8.0$ \\
T3: T2 + spurious link  & 16.0 & 273.5 & 289.5 &     $-16.0$ \\
\bottomrule
\end{tabular}
\end{center}

\textit{Bits saved = MDL(T0) $-$ MDL(theory).}

At regularity $= 1.0$, T2 saves 44.9 bits. T3 saves 36.9 bits---
exactly 8 bits less than T2. The spurious link costs exactly its
description length and changes $L(D \mid H)$ not at all. This is
Occam's razor in its precise form: \emph{a theory that is correct but
unnecessarily complex loses to the minimal correct theory by exactly
the description length of the excess.}

At regularity $= 0.5$, T2 barely saves 1.9 bits. The signal is real
but weak; T3 cannot recover the extra 8-bit cost. At regularity $=
0.0$, T0 wins; T1 and T2 tie at exactly $-8.0$ bits; T3 pays $-16.0$
bits. The framework penalizes each spurious link by exactly its
description length.

The regularity crossover is the MDL phase transition of
Section~\ref{sec:connections} instantiated at the bit level. As
regularity increases from 0 to 1, the optimal theory switches
discontinuously from T0 to T2. The switch is not gradual: T2 goes
from $-8.0$ bits (noise) to $+1.9$ bits (weak signal) to $+44.9$ bits
(perfect regularity). The optimal hypothesis class jumps when the
theory's amortized savings first exceed its description length cost.
This is exactly what Kuhn observed in the history of science and what
Lakatos formalized as the distinction between progressive and
degenerating research programmes. Here it is a measurable threshold
in bits.